\section{Theory and Hypotheses}

\subsection{Core Theoretical Claim}

The project starts from a simple but consequential proposition: humans learn not only whether AI is good or bad, but where AI belongs in a decision system. This means that the relevant object of learning is the collaboration architecture itself. We define a human--AI collaboration architecture as the role allocation structure that specifies who decides, who recommends, who confirms, and who can override.

\subsection{Architecture Attraction}

Each architecture carries an ``attraction'' value: a subjective expectation of its desirability given prior feedback about accuracy, error source, correction opportunities, time cost, review burden, and perceived risk. Repeated experience updates these attractions. When the attractions of some architectures rise relative to others, later choices increasingly concentrate on them. Convergence occurs when this choice distribution stabilizes.

\subsection{Asymmetric Updating}

The key mechanism is asymmetric updating. If AI errors and human errors are learned symmetrically, then under equal long-run performance, early sequence differences should wash out over time. But if AI errors trigger stronger negative updates than comparable human errors, then early AI failures may disproportionately suppress later willingness to choose AI-led architectures. This introduces a self-reinforcing dynamic: lower authorization creates less later exposure to potentially trust-repairing AI performance, which slows recovery and increases the chance of long-run lock-in.

\subsection{Path Dependence and Behavioral Steady States}

The project does not claim a Nash equilibrium in the strict game-theoretic sense because the AI system is not modeled as an independently strategic actor. Instead, the paper uses the safer and more behaviorally appropriate language of behavioral steady states, learning equilibrium, and architecture convergence. The central concern is whether a repeated learning process stabilizes in a suboptimal authorization regime.

\subsection{Hypotheses}

\textbf{H1: Asymmetric negative learning.} When AI and human decision makers exhibit equivalent objective long-run performance, AI errors will trigger stronger negative updating of collaboration architectures than comparable human errors.

\textbf{H2: Path-dependent convergence.} Early AI errors will increase the likelihood that participants converge toward lower-AI-autonomy architectures.

\textbf{H3: Recovery asymmetry.} The positive effect of later correct AI performance on restoring AI authorization will be weaker than the negative effect of early AI errors on reducing AI authorization.

\textbf{H4: Human-error tolerance.} Early human errors will not symmetrically push participants toward substantially higher-AI-autonomy architectures.

\textbf{H5: Suboptimal steady states.} Even when AI and human long-run objective performance is equivalent, participants may settle into collaboration architectures that are more supervision-heavy than the objectively most efficient configuration.
